\subsection{Local character fields and quadratic codegree saturation}
\label{sec:r3-local-fields}

The next criterion strengthens the signed Casimir bounds by determining
the dimension field before enumerating possible formal codegrees.
Write $A$ for the signed algebra supplied by
Lemma~\ref{lem:dimension-field-containment}, $K$ for its dimension
character field, and $\Delta$ for its dimension codegree. Thus every
original or signed formal codegree $f$ lies in $K$, and
$\Delta/f$ is an algebraic integer. In particular this statement
concerns the sphericalization of an arbitrary categorification; it
does not assume a pivotal structure on that category.

\begin{proposition}[Two local ways to isolate the dimension field]
\label{prop:r3-local-isolation}
Let $R$ be commutative and self-dual of rank six, and suppose an
original formal codegree generates $E=\mathbb Q(\sqrt q)$, where
$q>1$ is squarefree. Either of the following conditions implies
$K=E$ and that the other four signed formal codegrees are rational
integers greater than one:
\begin{enumerate}
\item $q\equiv5\pmod8$, and $R/2R$ is a product of finite fields
with exactly one even-degree factor, of degree two;
\item $R/2R$ has three local factors: $\mathbb F_8$, a rank-two
local algebra, and $\mathbb F_2$.
\end{enumerate}
In the second case the conclusion uses cyclotomicity of characters
of a categorifiable fusion ring.
\end{proposition}
\begin{proof}
The dimension-field containment lemma gives $E\subset K$.
Because the regular matrices of $A$ are symmetric, $A\otimes\mathbb Q$
is a product of totally real number fields. Its primitive local
idempotents at two lift from $A/2A=R/2R$ to $A\otimes\mathbb Z_2$.
A finite free local factor with finite-field special fibre is finite
\'etale: its relative differentials vanish by Nakayama, and it is
flat. It is therefore the unramified extension with that residue
degree. Equivalently one may apply the classification of finite
\'etale algebras over a henselian local ring; see the Stacks Project,
\href{https://stacks.math.columbia.edu/tag/04GK}{Tag 04GK}.

In the first case $E\otimes\mathbb Q_2$ is the unramified quadratic
field. Every completion of a number field containing $E$ has even
degree over $\mathbb Q_2$. Since $K$ is a field quotient of
$A\otimes\mathbb Q$, its completions must use only the unique
even-degree local factor. That factor has degree two, so $K=E$.
If a codegree belonging to another character field $F$ were
irrational, it would generate $E$, because it already belongs to
$K=E$. Then $F$ would require another even-degree local factor,
which is impossible. All remaining codegrees are therefore rational.

In the second case the lifted rank-three factor has generic fibre
the unramified cubic extension of $\mathbb Q_2$. Let $F$ be the
global character field containing this local factor. A signed
character extends to a character of the sphericalization with
the sign object sent to $-1$. Hence its values are cyclotomic by
\cite[Theorem~8.51]{ENO}; in particular $F/\mathbb Q$ is abelian
and Galois. All its completions at two have the same local degree,
here three. Thus $[F:\mathbb Q]$ is three or six. Degree six would
require two degree-three local factors, whereas the remaining
lifted blocks have ranks two and one. Hence $F$ is cubic and uses
the entire rank-three block. The remaining rational algebra has
dimension three. The field $K$ contains the quadratic field $E$,
so it lies in this remainder and has degree two. Thus $K=E$.
The other character fields are $F$ and $\mathbb Q$; their
intersections with $E$ are rational. The remaining codegrees,
which lie in these intersections, are rational as claimed.

In both cases rational formal codegrees are algebraic integers.
They are positive, and the sum of reciprocals of the six signed
codegrees is one. Each is consequently greater than one.
\end{proof}

The second argument does not assert that the whole two-adic order
is \'etale. In particular a rank-two order with dual-number
reduction may have split generic fibre. Only its rank and its
separation from the cubic block are needed.

\begin{proposition}[Finite quadratic spectrum bound]
\label{prop:r3-quadratic-spectrum}
Suppose $K=E=\mathbb Q(\sqrt q)$ and the complementary four
codegrees are rational, as in Proposition~\ref{prop:r3-local-isolation}.
Assume $q\equiv1\pmod4$ and the original Frobenius--Perron
codegree has the form $D=c\sqrt q\,u$, where $c$ is a positive
integer and $u$ is an integral unit of norm $-1$.
Let $T=\operatorname{Tr}(\sum_iN_i^2)$.
Every categorification supplies integers $x,y$ and
$f_1\leq\cdots\leq f_4$ satisfying the following finite system:
\[
\begin{gathered}
 \Delta=\frac{c(qy+x\sqrt q)}2,\qquad
 x\equiv y\pmod2,\quad y>0,\quad x^2<qy^2,\quad cqy\leq T-5,\\
 n=qy^2-x^2,\qquad n\mid4qy^2,\qquad
 \sum_{i=1}^4\frac1{f_i}=1-\frac{4y}{cn}>0,\\
 f_i>1,\qquad
 f_i\mid\gcd\left(\frac{c(qy-x)}2,cx\right),\qquad
 cqy+\sum_{i=1}^4f_i\leq T.
\end{gathered}
\]
In addition $\Delta/f$ is integral for every original rational
codegree $f$, and the polynomial
\[
 \left(z^2-cqy\,z+\frac{c^2qn}{4}\right)
       \prod_{i=1}^4(z-f_i)
\]
is congruent modulo two to the original Casimir polynomial.
If the system is empty, $R$ is not categorifiable. Otherwise its
least total trace bounds the possible signed Casimir deficit.
\end{proposition}
\begin{proof}
Since $\Delta/D$ is integral,
$\theta=\Delta/(c\sqrt q)=u\Delta/D$ is integral. The integral
basis of $E$ gives $\theta=(x+y\sqrt q)/2$ with the stated
parity. Total positivity of $\Delta$ gives $y>0$ and
$x^2<qy^2$. Its two dimension characters contribute trace $cqy$;
the other four codegrees exceed one. The signed Casimir trace is
an integer at most $T$, giving the displayed finite bound.

The global dimension is a $d$-number by
\cite[Theorem~1.2]{Ocode}. Here
\[
 \operatorname{Tr}(\Delta)=cqy,\quad
 \operatorname{N}(\Delta)=\frac{c^2qn}{4},\quad
 \operatorname{Tr}(1/\Delta)=\frac{4y}{cn}.
\]
The quadratic $d$-number divisibility gives $n\mid4qy^2$,
and the reciprocal-codegree identity gives the reciprocal equation.
The integral-basis coordinates of $\Delta$ are
$c(qy-x)/2$ and $cx$. Thus $\Delta/f_i$ is integral precisely
when $f_i$ divides their greatest common divisor. Divisibility
by original rational codegrees follows separately from induction.
Finally the signed and original matrices are congruent modulo two,
so their Casimir characteristic polynomials have the same reduction.
Every enumeration range is finite. The case $x=0$ is retained:
$K$ can be quadratic even when $\Delta$ itself is rational.
\end{proof}

\begin{theorem}\label{thm:r3-arithmetic-exclusions}
The rings $\RR{8}{6524}$, $\RR{8}{6525}$ and $\RR{8}{6999}$
have no complex fusion categorifications.
\end{theorem}
\begin{proof}
For the first two rings, $q=149$, $c=9$,
\[
 D=\frac{6705+549\sqrt{149}}2,\qquad T=6741.
\]
Their reductions at two have field degrees $(2,3,1)$ and
$(2,1,1,1,1)$, respectively. Proposition~\ref{prop:r3-local-isolation}(1)
applies. The finite spectrum test has $1\leq y\leq5$ and leaves
only
\[
 y=5,\qquad x=\pm61,\qquad (f_1,f_2,f_3,f_4)=(9,9,9,9).
\]
Both traces equal $6741$.

For $\RR{8}{6999}$, $q=41$, $c=18$,
\[
 D=3690+576\sqrt{41},\qquad T=7416.
\]
The local factors are $\mathbb F_8$, the dual numbers of rank two,
and $\mathbb F_2$. Proposition~\ref{prop:r3-local-isolation}(2)
applies. The bound is $1\leq y\leq10$; the only survivors are
\[
 y=10,\qquad x=\pm64,\qquad (f_1,f_2,f_3,f_4)=(9,9,9,9),
\]
again with total trace $T$.

In all three cases the signed Casimir trace is the sum of the
squares of the signed structure constants. Since each absolute
value is bounded by its original coefficient, equality of traces
forces equality of every absolute value. The remaining unrestricted
sign-only associativity systems are exhausted exactly. Every
solution is a dual-compatible sign gauge of the original tensor.
After this gauge, the lift fusion rules are the original positive
ones, and the recorded dimension-character tests leave only the
Galois orbit of the positive character. Any categorification would
therefore have a pseudo-unitary Galois conjugate.

The pseudo-unitary possibility is excluded by the exact
\texttt{global\_codegree\_moment} certificates, checked independently
on the full input tensors. For the first two rings the obstruction
is visible already in the integral span of simple dimensions.
They are $1,u-1,u,u,u,u$, where
$u=(37+3\sqrt{149})/2$ and $u^2-37u+7=0$.
The codegree $9$ would require a central simple of dimension
\[
 D/9=(61u-11)/3\notin\mathbb Z+\mathbb Z u,
\]
which cannot be the dimension of an integral restriction to the
original category. This also gives a direct check on the final
induction obstruction in those two cases.
\end{proof}

The tensor-only generator and exact verifier are
\filename{software/round3_arithmetic.py}; all local algebras,
surviving spectra, sign gauges, and complete general certificates
are recorded in \filename{research/round3/arithmetic/}.
The independent enumeration in
\filename{research/round3/models/arithmetic_audit.py} reproduces
the first pair of saturated spectra without an optimizer or the
generator's recursive divisor enumeration.

\subsection{Noncommutative signed spectra}
\label{sec:r3-noncommutative}

The preceding field-isolation argument has a different replacement
when the signed quotient is forced to be noncommutative. This follows
either from reduction modulo two or from a separate exhaustive
exclusion of the commutative signed sector. It is essential in this
case to keep the representation multiplicities.

\begin{proposition}\label{prop:r3-noncommutative-spectrum}
Let $R$ have rank six. Suppose either that its reduction modulo two
is noncommutative, or that every commutative signed quotient permitted
by sphericalization has been excluded by necessary conditions.
Suppose its two linear characters form a real quadratic orbit
over $E=\mathbb Q(\sqrt q)$, and let $D,D'$ be their formal
codegrees and $c$ that of its unique two-dimensional representation.
For any categorification there are $\Delta\in E$ and
$h\in\mathbb Z_{>2}$ satisfying
\[
 \frac1\Delta+\frac1{\Delta'}+\frac2h=1,
 \qquad \operatorname{Tr}_{E/\mathbb Q}(\Delta)+8h\leq T,
 \qquad T=\sum_{i,j,k}(N_{ij}^k)^2.
\]
Both embeddings of $\Delta$ are positive, $\Delta$ is a
$d$-number, and
\[
 \Delta/D,\quad \Delta/D',\quad \Delta/c,\quad \Delta/h
 \quad\hbox{belong to }\mathcal O_E.
\]
For fixed $R$ these conditions are a finite exact search.
\end{proposition}
\begin{proof}
The negative-sign quotient of sphericalization is now an oriented
signed algebra $A$. Its signed basis is orthonormal for the
coefficient-of-unit inner product, and duality gives adjoints of
left multiplication. Thus its regular representation is a finite
dimensional $*$-algebra with positive inner product and is
semisimple. In the first case parity preserves a nonzero commutator;
in the second case the commutative alternative is excluded separately.
In either case dimension six and the existence of a linear character give
\[
 A\otimes\mathbb C\cong\mathbb C\oplus\mathbb C\oplus M_2(\mathbb C).
\]
There are exactly two linear characters. The spherical dimension
character is real and dual invariant, and its field $K$ therefore
has degree at most two. Induction applied to the original characters
still gives $\Delta/f\in\mathcal O_K$. In particular $D\in K$,
so $K=E$. The signed linear codegrees are $\Delta,\Delta'$.

The unique two-dimensional representation is preserved by Galois
action. Its formal codegree $h$ is consequently rational and an
algebraic integer, hence an integer. The formal-codegree reciprocal
identity weights it by its representation dimension two. On the
four-dimensional regular matrix block the Casimir has eigenvalue
$2h$, so that block contributes $8h$ to the regular trace.
This proves both displayed identities, with the trace inequality
following from the coefficient bounds. Positivity of the two linear
codegrees and the reciprocal equation imply $h>2$.
Induction from the negative-sign character gives $\Delta/h$;
induction from the original characters gives the other quotients.
The usual global-dimension $d$-number theorem applies to $\Delta$.

For example, when $q\equiv1\pmod4$, write
$\Delta=(a+b\sqrt q)/2$, with $a\equiv b\pmod2$.
Then $0<a\leq T-24$ and $qb^2<a^2$. These bounds are finite.
For each pair $(a,b)$ the reciprocal equation determines
\[
 h=\frac{2}{1-a/\operatorname{N}(\Delta)}.
\]
Testing integrality, positivity, the four quotients, and the
quadratic $d$-number condition completes the exact necessary search.
\end{proof}

\begin{theorem}\label{thm:r3-noncommutative-exclusion}
The ring $\RR{6}{170}$ has no complex fusion categorification.
\end{theorem}
\begin{proof}
Its original codegrees are $148\pm24\sqrt{37}$ and $4$, with
representation dimensions $1,1,2$, and $T=328$.
The preceding finite system leaves exactly
\[
 \Delta=148\pm24\sqrt{37},\qquad h=4,
 \qquad \operatorname{Tr}(\Delta)+8h=328.
\]
Thus every signed coefficient has its original absolute value.
The complete oriented sign-only associativity enumeration has
$1631$ search nodes and $16$ solutions; all sixteen are
dual-compatible sign gauges of $R$. No commutativity or indicator
normalization is imposed. Under each gauge the two possible
dimension characters are Galois conjugates of the positive one,
so every categorification would have a pseudo-unitary Galois
conjugate.

Multiplicity-weighted central induction has $146$ candidate
restriction columns for the nontrivial linear character and $25$
for the two-dimensional representation. Their complete compatibility
search yields nineteen low frames. Every frame violates an ordinary
exact Gram or moment condition. The pseudo-unitary possibility is
therefore impossible, completing the general exclusion.
\end{proof}

The implementation is \filename{software/round3_noncommutative.py},
with arithmetic spectra, all oriented sign tables, the nineteen
induction branches, and replays in
\filename{research/round3/noncommutative/}.
For rings with commutative reduction modulo two, the first hypothesis
of Proposition~\ref{prop:r3-noncommutative-spectrum} is unavailable.
Their entire commutative signed sector must be treated before the
noncommutative spectrum calculation supplies a general reduction.

\begin{theorem}\label{thm:r3-noncommutative-sector-exclusions}
The rings $\RR{7}{200}$ and $\RR{8}{491}$ have no complex fusion
categorifications.
\end{theorem}
\begin{proof}
Both reductions modulo two are commutative. First impose commutativity
on the unrestricted oriented signed coefficient problem, identifying
opposite coefficient variables and intersecting their original parity
intervals. If an opposite original coefficient is zero, its identified
signed coefficient is forced to zero. This exhausts all commutative
signed quotients without an indicator normalization.

The complete enumerations have respectively $176$ and $208$ signed
tables, forming $11$ and $13$ dual-compatible sign-gauge orbits.
Their character algebras have respectively $27$ and $53$ components.
Each component has a zero dimension, unequal dimensions on dual
objects, or an exact norm-divisibility obstruction to induction.
The latter uses the original formal codegrees, whose quadratic
polynomial is $z^2-1560z+1872$, and the matrix-block codegree $12$.
The regular Casimir eigenvalue of that block is $24$; it must not
replace the formal codegree in the induction test. The recorded
norm obstructions occur at primes $3$ or $13$. Consequently no
categorical signed quotient is commutative.

Proposition~\ref{prop:r3-noncommutative-spectrum} now applies.
For both rings its finite system has exactly
\[
 \Delta=780\pm216\sqrt{13},\qquad h=12,\qquad
 \operatorname{Tr}(\Delta)+8h=1656=T.
\]
The resulting sign-only searches have $4927$ and $6463$ nodes,
respectively, and sixteen solutions each. Every solution is a
dual-compatible gauge of the original tensor. Thus any
categorification has a pseudo-unitary Galois conjugate. In each
ring the multiplicity-weighted induction problem has ten complete
low frames, all excluded by exact ordinary Gram or moment conditions.
This excludes the pseudo-unitary possibility and proves the claim.
\end{proof}

The separate commutative-sector verifier is
\filename{software/round3_commutative_sector.py}; its complete
certificate is included in each general certificate in
\filename{research/round3/noncommutative/}.
Both general replays also reconstruct the noncommutative arithmetic
and induction calculations. Removing the sector certificate is
explicitly rejected by the verifier.

\paragraph{Scope and controls.}
Each detector reconstructs its hypotheses from the full tensor.
The new family and arithmetic criteria are reapplied to the complete
census through multiplicity eight; exclusions of already resolved
rings are counted as confirmations. Exact replay disables numerical
optimizers and includes altered-certificate rejection and changes
of basis. All $54$ positively realized rings are controls, with
ineligibility distinguished from survival of an applicable system.
The C4 family includes the positive parameter-zero member in its
applicability scan, although the present irrational-dimension trace
engine does not apply to it. Its positive parameter-two member
survives the necessary system. These controls check scope; they do
not turn a surviving necessary condition into a categorical model.

\subsection{Dimension-block Fourier bounds and saturated twist support}
\label{subsec:dimension-block-support}
Let $K_1=\QQ(d_i)$, $K_0=\QQ(D)$ and $e=[K_1:K_0]$ for a spherical
dimension function. For a character $\chi$ of
$\operatorname{Gal}(K_0/\QQ)$, let $f_\chi$ be its primitive Dirichlet
conductor. Schopieray's dimension-block theorem
\cite[Theorem~5.1]{SchDimensionBlock} gives the additional necessary bound
\begin{equation}\label{eq:dimension-block-fourier-bound}
 \left|\sum_{\sigma\in\operatorname{Gal}(K_0/\QQ)}
       \frac{\chi(\sigma)}{\sigma(D)}\right|^2
 \le \frac1{e^2 f_\chi}.
\end{equation}
Its equality case supplies a restriction on individual twists, rather
than only on their averages. Here $S$ denotes the center's modular
$S$-matrix normalized by $D$, and $T$ is its matrix of actual twists.

\begin{proposition}[Saturated conductor support]
\label{prop:saturated-conductor-support}
Suppose equality holds in \eqref{eq:dimension-block-fourier-bound}.
Inflate $\chi$ to $G_1=\operatorname{Gal}(K_1/\QQ)$, and let
$w_\chi=\sum_{\sigma\in G_1}\chi(\sigma)e_{A_\sigma}$, where
$A_\sigma$ is the low central summand indexed by the dimension
character $\sigma(d)$. The cyclic subrepresentation generated by
$w_\chi$ in the canonical modular representation of $Z(\mathcal C)$
factors through $\mathrm{SL}_2(\mathbb Z/f_\chi\mathbb Z)$.
Consequently, for a central simple with restriction column $v$,
\begin{equation}\label{eq:saturated-column-support}
 L_\chi(v):=\sum_{\sigma\in G_1}\chi(\sigma)
       \sigma\!\left(\frac{\sum_i v_i d_i}{D}\right)\ne0
 \quad\Longrightarrow\quad \theta_Z^{f_\chi}=1.
\end{equation}
If several saturated characters give nonzero coefficients, the twist
order divides the greatest common divisor of their conductors.
\end{proposition}
\begin{proof}
We extract the equality information from the local compression
calculation in \cite[Lemma~4.1]{SchDimensionBlock}. Let
$\Gamma=\mathrm{SL}_2(\mathbb Z/\ell^a\mathbb Z)$, let $U$ be its
upper unipotent subgroup, and let $P$ project onto $U$-fixed vectors.
For a nontrivial torus character $\alpha$ of conductor $\ell^c$, the
universal permutation representation on primitive columns realizes
$\operatorname{Ind}_U^\Gamma\mathbf1$. Its $\alpha$-space of
$U$-invariants splits into the one-dimensional big-cell space and
the small-cell space. The compression $PsP$ maps the big-cell line
to the small-cell space and maps the latter back to a line. Both
nonzero singular values are $\ell^{-c/2}$. Hence every vector
attaining this norm belongs to the span of the big-cell vector and
the adjoint image of the opposite-character big-cell vector.

Writing a primitive column as $(p,q)$, these two vectors are, up to
nonzero scalar factors,
\[
 \alpha(q)[q\text{ is a unit}],\qquad
 \alpha(p)[q\equiv0\pmod{\ell^c}].
\]
Both are pullbacks of functions on primitive columns modulo $\ell^c$.
Their entire $\Gamma$-cyclic spans therefore factor through reduction
modulo $\ell^c$. Every irreducible representation with $U$-invariants
embeds equivariantly and isometrically, up to a scalar, into this
permutation representation. Thus an irreducible constituent containing
a nonzero norm-attaining vector itself factors through that reduction.

For the trivial local character the bound is $1$. Equality implies
that $sw$ is $U$-fixed as well as $w$. The vector $w$ is therefore fixed
by both upper and lower unipotents. These generate $\Gamma$, so its
cyclic representation is trivial. Decompose a general unitary
representation orthogonally into irreducibles, and then use the Chinese
remainder tensor decomposition at the prime divisors of its level.
Equality in the product norm bound forces every occupied local factor
to attain its own bound. The preceding conclusions show that a
norm-attaining vector generates a representation of level dividing
$f_\chi$.

The proof of the dimension-block bound estimates the matrix coefficient
$\langle w_{\bar\chi},PSPw_\chi\rangle$ by
$f_\chi^{-1/2}\|w_\chi\|^2$. Equality in the bound forces equality
in this compression-norm estimate and in Cauchy--Schwarz. The preceding
argument applies to $w_\chi$. In particular
$(T^{f_\chi}-I)Sw_\chi=0$. The known dimension-character rows give
$S_{Z,A_\sigma}=\sigma(d_Z/D)$, so the $Z$-coordinate of $Sw_\chi$
is precisely $L_\chi(v)$. This proves
\eqref{eq:saturated-column-support}. The greatest-common-divisor
assertion follows by intersecting root-of-unity groups.
\end{proof}

The exact audit \filename{research/round5/unification/dimension_block_bound.py}
enumerates every dimension-field character for the six unresolved rings.
All satisfy the bound. Their nontrivial equality conductors are
$20$ for $R_{6,4,22}$, $16$ for $R_{6,4,91}$, $15$ for $R_{6,4,629}$,
$12$ for $R_{6,5,1204}$, and $12,24$ for $R_{6,8,1660}$;
$R_{6,4,155}$ has strict inequalities. The receipt records all Fourier
values, including trivial characters, in exact cyclotomic coordinates.
The reusable implementation \filename{software/dimension_block_support.py}
computes \eqref{eq:saturated-column-support} from the dimensions and each
restriction column. When the forced roots lie in the retained field,
their complete power tuples filter the projected twist types exactly.

\begin{proposition}[Universal moments of a saturated real character]
\label{prop:saturated-modular-moments}
In Proposition~\ref{prop:saturated-conductor-support}, suppose $\chi$
is real and put $f=f_\chi$,
$\lambda=\sum_{\sigma\in G_1}\chi(\sigma)/\sigma(D)$.
Let $M=\operatorname{Ind}_{B}^{\mathrm{SL}_2(\mathbb Z/f\mathbb Z)}\chi$,
where $B$ is the upper triangular subgroup and its character is trivial
on unipotents and equals $\chi$ on the diagonal torus. Write $S_f,T_f$
for the standard unitary operators on $M$, and $P_\chi$ for projection
onto its unipotent-fixed, torus-$\chi$ space. Let the columns of $K$
span the subspace on which $P_\chi S_fP_\chi$ has eigenvalue $\lambda$.
For every exact matrix identity
\[
 \sum_{n=0}^{f-1}c_n K^*S_f^*T_f^nS_fK=0,
\]
the following necessary moment equation holds:
\begin{equation}\label{eq:saturated-modular-moments}
 \sum_{n=0}^{f-1}c_n\sum_Z|L_\chi(v_Z)|^2\theta_Z^n=0.
\end{equation}
\end{proposition}
\begin{proof}
By the preceding proposition, the cyclic representation generated by
$w_\chi$ factors through this finite group. It embeds equivariantly and
isometrically into $M$. Indeed, in an isotypic summand
$\pi\otimes\mathbb C^m$, the generating vector lies in
$\pi^{B,\chi}\otimes\mathbb C^m$; cyclicity forces
$m\le\dim\pi^{B,\chi}$. Frobenius reciprocity gives exactly this latter
multiplicity in $M$. Complete reducibility and invariant inner products
make the embedding isometric.

Since $\chi$ is real, equality in Cauchy--Schwarz in the preceding
proof says $P_\chi Sw_\chi=\lambda w_\chi$. Thus its embedded vector
lies in the column space of $K$. Applying the displayed matrix identity
to that vector gives
$\sum_n c_n\langle Sw_\chi,T^nSw_\chi\rangle=0$.
The simple-coordinate expansion of this inner product is exactly
\eqref{eq:saturated-modular-moments}.
\end{proof}
The signed permutation model on the projective line, its projectors,
equality eigenspace and complete matrix-relation kernel are constructed
by \filename{research/round5/unification/equality_modrep.py}.
This uses the full equality eigenspace, so it does not select an
unproved irreducible constituent or a particular vector within it.
There is also a direct spectral support test. If $E_r$ is the orthogonal
projector onto the $\zeta_f^r$-eigenspace of $T_f$, compute
\[
 K^*S_f^*E_rS_fK=
 \frac1f\sum_{n=0}^{f-1}\zeta_f^{-rn}K^*S_f^*T_f^nS_fK.
\]
If this matrix is zero, positivity implies $E_rS_fK=0$.
Thus $L_\chi(v_Z)\ne0$ forbids $\theta_Z=\zeta_f^r$.
This exact discrete Fourier calculation can remove phase variables
before the moment equations are solved.

\begin{corollary}[Different saturated conductors are orthogonal]
\label{cor:saturated-cross-moments}
Let $\chi,\psi$ be saturated characters with distinct primitive
conductors $f_\chi,f_\psi$. Then for every integer $n$,
\begin{equation}\label{eq:saturated-cross-moments}
 \sum_Z\overline{L_\chi(v_Z)}L_\psi(v_Z)\theta_Z^n=0.
\end{equation}
Only the exponents $0\le n<\gcd(f_\chi,f_\psi)$ are needed.
\end{corollary}
\begin{proof}
Every irreducible constituent in the cyclic span of $w_\chi$ contains
a nonzero unipotent-fixed vector of torus character $\chi$. Its local
congruence level is at most the conductor exponent, by
Proposition~\ref{prop:saturated-conductor-support}, and cannot be smaller:
otherwise this torus character would factor through a smaller modulus.
The constituent therefore has exact conductor level $f_\chi$.
Constituents belonging to distinct conductors cannot be isomorphic.
Orthogonal complete reducibility of the unitary modular representation
now makes the two cyclic subspaces orthogonal. In particular
$\langle Sw_\chi,T^nSw_\psi\rangle=0$, which is the displayed formula.
Every nonzero summand has both Fourier coefficients nonzero, so its
twist order divides both conductors. The expression is therefore
periodic with period their greatest common divisor.
\end{proof}
